% Options for packages loaded elsewhere
% Options for packages loaded elsewhere
\PassOptionsToPackage{unicode}{hyperref}
\PassOptionsToPackage{hyphens}{url}
%
\documentclass[
  english,
  russian,
  12pt,
  a4paper,
  DIV=11,
  numbers=noendperiod]{scrreprt}
\usepackage{xcolor}
\usepackage{amsmath,amssymb}
\setcounter{secnumdepth}{5}
\usepackage{iftex}
\ifPDFTeX
  \usepackage[T1]{fontenc}
  \usepackage[utf8]{inputenc}
  \usepackage{textcomp} % provide euro and other symbols
\else % if luatex or xetex
  \usepackage{unicode-math} % this also loads fontspec
  \defaultfontfeatures{Scale=MatchLowercase}
  \defaultfontfeatures[\rmfamily]{Ligatures=TeX,Scale=1}
\fi
\usepackage{lmodern}
\ifPDFTeX\else
  % xetex/luatex font selection
\fi
% Use upquote if available, for straight quotes in verbatim environments
\IfFileExists{upquote.sty}{\usepackage{upquote}}{}
\IfFileExists{microtype.sty}{% use microtype if available
  \usepackage[]{microtype}
  \UseMicrotypeSet[protrusion]{basicmath} % disable protrusion for tt fonts
}{}
\usepackage{setspace}
% Make \paragraph and \subparagraph free-standing
\makeatletter
\ifx\paragraph\undefined\else
  \let\oldparagraph\paragraph
  \renewcommand{\paragraph}{
    \@ifstar
      \xxxParagraphStar
      \xxxParagraphNoStar
  }
  \newcommand{\xxxParagraphStar}[1]{\oldparagraph*{#1}\mbox{}}
  \newcommand{\xxxParagraphNoStar}[1]{\oldparagraph{#1}\mbox{}}
\fi
\ifx\subparagraph\undefined\else
  \let\oldsubparagraph\subparagraph
  \renewcommand{\subparagraph}{
    \@ifstar
      \xxxSubParagraphStar
      \xxxSubParagraphNoStar
  }
  \newcommand{\xxxSubParagraphStar}[1]{\oldsubparagraph*{#1}\mbox{}}
  \newcommand{\xxxSubParagraphNoStar}[1]{\oldsubparagraph{#1}\mbox{}}
\fi
\makeatother


\usepackage{longtable,booktabs,array}
\usepackage{calc} % for calculating minipage widths
% Correct order of tables after \paragraph or \subparagraph
\usepackage{etoolbox}
\makeatletter
\patchcmd\longtable{\par}{\if@noskipsec\mbox{}\fi\par}{}{}
\makeatother
% Allow footnotes in longtable head/foot
\IfFileExists{footnotehyper.sty}{\usepackage{footnotehyper}}{\usepackage{footnote}}
\makesavenoteenv{longtable}
\usepackage{graphicx}
\makeatletter
\newsavebox\pandoc@box
\newcommand*\pandocbounded[1]{% scales image to fit in text height/width
  \sbox\pandoc@box{#1}%
  \Gscale@div\@tempa{\textheight}{\dimexpr\ht\pandoc@box+\dp\pandoc@box\relax}%
  \Gscale@div\@tempb{\linewidth}{\wd\pandoc@box}%
  \ifdim\@tempb\p@<\@tempa\p@\let\@tempa\@tempb\fi% select the smaller of both
  \ifdim\@tempa\p@<\p@\scalebox{\@tempa}{\usebox\pandoc@box}%
  \else\usebox{\pandoc@box}%
  \fi%
}
% Set default figure placement to htbp
\def\fps@figure{htbp}
\makeatother



\ifLuaTeX
\usepackage[bidi=basic,provide=*]{babel}
\else
\usepackage[bidi=default,provide=*]{babel}
\fi
% get rid of language-specific shorthands (see #6817):
\let\LanguageShortHands\languageshorthands
\def\languageshorthands#1{}


\setlength{\emergencystretch}{3em} % prevent overfull lines

\providecommand{\tightlist}{%
  \setlength{\itemsep}{0pt}\setlength{\parskip}{0pt}}



 
\usepackage[style=gost-numeric,backend=biber,langhook=extras,autolang=other*]{biblatex}
\addbibresource{bib/cite.bib}

\usepackage[]{csquotes}

\usepackage{indentfirst}
\usepackage{float}
\floatplacement{figure}{H}
\IfFileExists{plex-otf.sty}{
  %% Full TeXlive
  \usepackage[
    % math,
    RM={Scale=0.94},SS={Scale=0.94},SScon={Scale=0.94},TT={Scale=MatchLowercase,FakeStretch=0.9},
    DefaultFeatures={Ligatures=Common}
  ]{plex-otf}
  % \usepackage[Scale=MatchUppercase]{juliamono}
}{
  %% TinyTeX
  \usepackage{libertine}
}
\KOMAoption{captions}{tableheading}
\makeatletter
\@ifpackageloaded{caption}{}{\usepackage{caption}}
\AtBeginDocument{%
\ifdefined\contentsname
  \renewcommand*\contentsname{Содержание}
\else
  \newcommand\contentsname{Содержание}
\fi
\ifdefined\listfigurename
  \renewcommand*\listfigurename{Список иллюстраций}
\else
  \newcommand\listfigurename{Список иллюстраций}
\fi
\ifdefined\listtablename
  \renewcommand*\listtablename{Список таблиц}
\else
  \newcommand\listtablename{Список таблиц}
\fi
\ifdefined\figurename
  \renewcommand*\figurename{Рисунок}
\else
  \newcommand\figurename{Рисунок}
\fi
\ifdefined\tablename
  \renewcommand*\tablename{Таблица}
\else
  \newcommand\tablename{Таблица}
\fi
}
\@ifpackageloaded{float}{}{\usepackage{float}}
\floatstyle{ruled}
\@ifundefined{c@chapter}{\newfloat{codelisting}{h}{lop}}{\newfloat{codelisting}{h}{lop}[chapter]}
\floatname{codelisting}{Список}
\newcommand*\listoflistings{\listof{codelisting}{Листинги}}
\makeatother
\makeatletter
\makeatother
\makeatletter
\@ifpackageloaded{caption}{}{\usepackage{caption}}
\@ifpackageloaded{subcaption}{}{\usepackage{subcaption}}
\makeatother
\usepackage{bookmark}
\IfFileExists{xurl.sty}{\usepackage{xurl}}{} % add URL line breaks if available
\urlstyle{same}
\hypersetup{
  pdftitle={Лабораторная работа №2},
  pdfauthor={Осипов Никита Александрович},
  pdflang={ru-RU},
  hidelinks,
  pdfcreator={LaTeX via pandoc}}


\title{Лабораторная работа №2}
\usepackage{etoolbox}
\makeatletter
\providecommand{\subtitle}[1]{% add subtitle to \maketitle
  \apptocmd{\@title}{\par {\large #1 \par}}{}{}
}
\makeatother
\subtitle{Дискреционное разграничение прав в Linux. Основные атрибуты}
\author{Осипов Никита Александрович}
\date{}
\begin{document}
\maketitle

\renewcommand*\contentsname{Содержание}
{
\setcounter{tocdepth}{1}
\tableofcontents
}
\listoffigures
\listoftables

\setstretch{1.5}
\chapter{Цель
работы}\label{ux446ux435ux43bux44c-ux440ux430ux431ux43eux442ux44b}

Получение практических навыков работы в консоли с атрибутами файлов,
закрепление теоретических основ дискреционного разграничения доступа в
современных системах с открытым кодом на базе ОС Linux.

\chapter{Задание}\label{ux437ux430ux434ux430ux43dux438ux435}

Постарайтесь последовательно выполнить все пункты, занося ваши ответы на
поставленные вопросы и замечания в отчёт. 1. В установленной при
выполнении предыдущей лабораторной работы операционной системе создайте
учётную запись пользователя guest (использую учётную запись
администратора): useradd guest 2. Задайте пароль для пользователя guest
(использую учётную запись администратора): passwd guest 3. Войдите в
систему от имени пользователя guest. 4. Определите директорию, в которой
вы находитесь, командой pwd. Сравните её с приглашением командной
строки. Определите, является ли она вашей домашней директорией? Если
нет, зайдите в домашнюю директорию. 5. Уточните имя вашего пользователя
командой whoami. 6. Уточните имя вашего пользователя, его группу, а
также группы, куда входит пользователь, командой id. Выведенные значения
uid, gid и др. запомните. Сравните вывод id с выводом команды groups. 7.
Сравните полученную информацию об имени пользователя с данными,
выводимыми в приглашении командной строки. 8. Просмотрите файл
/etc/passwd командой cat /etc/passwd Найдите в нём свою учётную запись.
Определите uid пользователя. Определите gid пользователя. Сравните
найденные значения с полученными в предыдущих пунктах. Замечание: в
случае, когда вывод команды не умещается на одном экране монитора,
используйте прокрутку вверх--вниз (удерживая клавишу shift, нажимайте
page up и page down) либо программу grep качестве фильтра для вывода
только строк, содержащих определённые буквенные сочетания: cat
/etc/passwd \textbar{} grep guest 9. Определите существующие в системе
директории командой ls -l /home/ Удалось ли вам получить список
поддиректорий директории /home? Какие права установлены на директориях?
10. Проверьте, какие расширенные атрибуты установлены на поддиректориях,
находящихся в директории /home, командой: lsattr /home Удалось ли вам
увидеть расширенные атрибуты директории? Удалось ли вам увидеть
расширенные атрибуты директорий других пользователей? 11. Создайте в
домашней директории поддиректорию dir1 командой mkdir dir1 Определите
командами ls -l и lsattr, какие права доступа и расширенные атрибуты
были выставлены на директорию dir1. 12. Снимите с директории dir1 все
атрибуты командой chmod 000 dir1 и проверьте с её помощью правильность
выполнения команды ls -l 13. Попытайтесь создать в директории dir1 файл
file1 командой echo \enquote{test} \textgreater{} /home/guest/dir1/file1
Объясните, почему вы получили отказ в выполнении операции по созданию
файла? Оцените, как сообщение об ошибке отразилось на создании файла?
Проверьте командой ls -l /home/guest/dir1 действительно ли файл file1 не
находится внутри директории dir1. 14. Заполните таблицу «Установленные
права и разрешённые действия» (см. табл. 2.1), выполняя действия от
имени владельца директории (файлов), определив опытным путём, какие
операции разрешены, а какие нет. Если операция разрешена, занесите в
таблицу знак «+», если не разрешена, знак «-». Замечание 1: при
заполнении табл. 2.1 рассматриваются не все атрибуты файлов и
директорий, а лишь «первые три»: г, w, х, для «владельца». Остальные
атрибуты также важны (особенно при использовании доступа от имени разных
пользователей, входящих в те или иные группы). Проверка всех атрибутов
при всех условиях значительно увеличила бы таблицу: так 9 атрибутов на
директорию и 9 атрибутов на файл дают 218 строк без учёта дополнительных
атрибутов, плюс таблица была бы расширена по количеству столбцов, так
как все приведённые операции необходимо было бы повторить ещё как
минимум для двух пользователей: входящего в группу владельца файла и не
входящего в неё. После полного заполнения табл. 2.1 и анализа полученных
данных нам удалось бы выяснить, что заполнение её в таком виде излишне.
Можно разделить большую таблицу на несколько малых независимых таблиц. В
данном примере предлагается рассмотреть 3 + 3 атрибута, т.е. 2 6 = 64
варианта

\chapter{Выполнение лабораторной
работы}\label{ux432ux44bux43fux43eux43bux43dux435ux43dux438ux435-ux43bux430ux431ux43eux440ux430ux442ux43eux440ux43dux43eux439-ux440ux430ux431ux43eux442ux44b}

\begin{enumerate}
\def\labelenumi{\arabic{enumi}.}
\tightlist
\item
  Создаем гостевую учетную запись и задаем ей пароль через командую
  строку.
\end{enumerate}

\begin{figure}[H]

{\centering \pandocbounded{\includegraphics[keepaspectratio]{image/1.png}}

}

\caption{Создание гостевой учетной записи}

\end{figure}%

\begin{enumerate}
\def\labelenumi{\arabic{enumi}.}
\setcounter{enumi}{1}
\tightlist
\item
  Заходим в гостевую запись и проверяем базовые параметры
\end{enumerate}

\begin{figure}[H]

{\centering \pandocbounded{\includegraphics[keepaspectratio]{image/2.png}}

}

\caption{Проверка учетной записи}

\end{figure}%

\begin{enumerate}
\def\labelenumi{\arabic{enumi}.}
\setcounter{enumi}{2}
\tightlist
\item
  Создаем папку и проверяем доступ к ней
\end{enumerate}

\begin{figure}[H]

{\centering \pandocbounded{\includegraphics[keepaspectratio]{image/3.png}}

}

\caption{Проверка}

\end{figure}%

\chapter{Выводы}\label{ux432ux44bux432ux43eux434ux44b}

Я получил практические навыки работы в консоли с атрибутами файлов,
закрепил теоретические основы дискреционного разграничения доступа в
современных системах с открытым кодом на базе ОС Linux.

\chapter*{Список
литературы}\label{ux441ux43fux438ux441ux43eux43a-ux43bux438ux442ux435ux440ux430ux442ux443ux440ux44b}
\addcontentsline{toc}{chapter}{Список литературы}

\printbibliography[heading=none]





\end{document}
